\expandafter\def\csname HistoricalNoteText363\endcsname{This paper contains the substance of the lecture I gave at the public meeting of the Royal Society of Sciences on 19 September 1837.}
\clearpage
\fancyhead[L]{\small ON A NEW INSTRUMENT.}
\SourceLeaf{363}
\pdfbookmark[0]{On a new instrument. 1837}{instrument-1837}
\begin{center}
{\large ON A NEW INSTRUMENT INTENDED PRIMARILY FOR DIRECTLY OBSERVING CHANGES IN THE INTENSITY OF THE HORIZONTAL COMPONENT OF THE EARTH’S MAGNETISM\HistoricalNote{363}.}
\par\medskip
Resultate aus den Beobachtungen des magnetischen Vereins. 1837. I.
\end{center}

As is well known, a complete determination of the Earth’s magnetism at a given place requires a system of \textit{three} elements. The usual choice is declination, inclination and intensity. Although this is the simplest choice to conceive of, a different combination is also possible and may be preferable in some respects. In particular, for practical as well as theoretical purposes it is far more advantageous to consider the horizontal component of the Earth’s magnetic force on its own and represent it by two elements, its direction (declination) and its intensity. If either the intensity of the vertical force or the inclination of the whole force is added as a third element, the intensity of the whole force follows directly whenever it is wanted.

As for the two elements of the horizontal component of the Earth’s magnetism, the only ones at issue here, the magnetometer that has been in use for five years has fully \SourceLeaf{364}solved every problem that arises in connection with declination. It can determine its absolute value and also track its regular and chance variations from year to year, from month to month, from hour to hour and even from one minute to the next, with a reliability, convenience and precision that leave nothing to be desired.

The same instrument also serves to determine the horizontal intensity of the Earth’s magnetism in absolute measure. Indeed, as is well known, that very task first prompted the development of the magnetometer. Yet the instrument falls well short of solving the problem completely in \textit{every} respect.

To show clearly what still remains to be desired, I must first recall that using the magnetometer to determine magnetic intensity involves a combination of \textit{several} operations, one of which is observing the duration of a needle’s swing. This by its very nature takes a considerable time, since the number of swings from which the duration of a single one must be inferred cannot be too small. If the magnetic intensity is constant throughout such an operation, the measured swing duration does indeed correspond to \textit{that} value of the intensity. But if the intensity varied in the meantime, the swing duration will correspond only to \textit{the mean value} of the intensity over that interval. In this way we are left entirely in the dark as to whether there were any changes in the magnetic intensity \textit{during} that time or what those changes were. The instrument can therefore give only average values over certain intervals, without giving a complete and faithful account of the course of the changes within them. If we tried to come closer to such a record by choosing shorter intervals or always basing our results on only a few swings, the results would lose too much precision and reliability. We would run the risk of taking mere errors of observation for anomalies in the intensity.

Experience over the past few years has shown how interesting the rapidly changing disturbances in the Earth’s magnetic force can be even when seen from just one side, in the declination. The more striking this became, the more keenly one wanted a means of tracking and measuring, with the same ease, reliability and precision, the effects that these disturbances \SourceLeaf{365}undoubtedly have on the intensity.,

As I have just explained, the means of observation used so far are unsuitable for \textit{this} purpose because they rely on timing the swings, and timing them by its nature takes too long on each occasion. Here the duration of a needle’s swing serves only to determine indirectly the torque that the Earth’s magnetic force exerts on the needle when it lies outside the magnetic meridian. If this torque can therefore be determined precisely by a direct method without observing any swings, and if its changes can be measured reliably, precisely and quickly, the main part of our problem will be solved. The method I used rests on the following principle.

The conditions for equilibrium of a body of arbitrary shape suspended from two threads, assuming for the moment that its parts are subject only to gravity and are firmly connected to one another, can be stated briefly. The vertical through the body’s centre of gravity and the straight lines traced by the threads must all lie in one plane and must either be parallel to one another or meet at a single point. Thus in any equilibrium position both threads and the centre of gravity lie in one vertical plane. To make the situation definite, suppose that the two threads are equally long, that their upper attachment points are at the same height and as far apart as the two lower ones, and finally that the two lower points form an isosceles triangle with the centre of gravity. Under these assumptions both threads hang vertically in equilibrium, and a third vertical line imagined midway between them passes through the body’s centre of gravity. If the body is turned out of this position about that last line, the two threads cease to hang vertically and cease to lie in a single plane, and the body rises slightly. It then tends to return to its former position under a torque that can be taken, to sufficient accuracy, as proportional to the sine of the displacement from rest. It is therefore greatest when the displacement is 90 degrees. This greatest torque is always tacitly \SourceLeaf{366}understood whenever torque is mentioned without qualification. It can also be regarded as a measure of a force that holds the body in its equilibrium position by virtue of the way it is suspended. For brevity I shall call this force the directive force due to the suspension. Its magnitude depends on 1) the length of the suspension threads, 2) the distance between them, 3) the weight of the body. It is inversely proportional to the length of the threads, directly proportional to the square of their separation and also directly proportional to the body’s weight. If the assumptions above do not hold exactly, the expression for the directive force is more complicated, and the reaction of the threads to twisting also calls for a small modification. There are, however, ways of determining the magnitude of the directive force experimentally with the greatest precision. If the body is left to itself after a smaller or larger displacement from equilibrium, it will oscillate with the greatest regularity. The midpoint of its swings coincides with this position, and their duration depends on the magnitude of the directive force and the moment of inertia of the body.

Suppose now that a horizontal bar magnet forms part of the suspended body. A second directive force comes into play, and what happens depends on how the two directive forces combine according to the familiar rules of statics. Three cases must be distinguished. The two positions in which the body would be in equilibrium under either force acting alone may coincide, be opposite, or make an angle with each other. It is easy to see that the difference between these three cases depends on the relationship between two angles. One is the angle the straight line through the two lower attachment points of the threads makes with the bar magnet. The other is the angle the straight line through the two upper attachment points makes with the magnetic meridian. Imagine the body in the equilibrium position determined by the suspension alone. In the first of our three cases, the bar magnet must lie in the magnetic meridian in its natural position (the north pole on the north side). In the second case it must lie in the meridian in the reversed position. In the third it must make an angle with the \SourceLeaf{367}magnetic meridian. For brevity I shall call these three possible positions of the bar magnet in the apparatus natural, reversed and transverse.

In the natural position, the effect of the Earth’s magnetism on the bar magnet leaves the equilibrium position corresponding to the suspension as it was. The apparatus is held there by a stronger force equal to the sum of the two directive forces.

In the second case, the reversed position, the apparatus is also in equilibrium in that position. Its equilibrium is stable only when the magnetic directive force is smaller than the directive force due to the suspension, and the apparatus is then held there only by a force equal to the difference between the two directive forces. If instead the magnetic directive force were the greater, the equilibrium would be unstable. Once displaced, the apparatus would fail to return. It would move farther and farther away, coming to rest only in the opposite position, where the bar magnet has its natural position in space while the suspension threads cross each other.

Finally, in the third case, where the two directive forces make an angle with each other, their opposition is resolved in an intermediate position. The bar magnet lies outside the meridian, and the straight line through the lower attachment points of the threads is no longer parallel to the line through the upper ones. Both the intermediate position and the force that holds the apparatus there follow the law of statics for the composition of two forces. It is now also clear that if the apparatus provides a way of measuring the angles between the three positions in question, the ratio of the two component directive forces can be calculated. The magnetic directive force can therefore be given in absolute measure if the directive force due to the suspension is known in absolute measure. Our problem is then solved. The most advantageous arrangement is to seat the bar magnet in relation to the other parts of the apparatus so that, in the resulting equilibrium position, it makes nearly a right angle with the magnetic meridian. The name transverse position is especially apt for this case. For one thing, the departure of the threads \SourceLeaf{368}from lying in a single plane is then greatest, so the result can be calculated with the greatest precision. For another, the small change in magnetic declination due to hourly or chance variations has no perceptible effect on the position. Conversely, every change in the intensity of the Earth’s magnetism affects the position directly and can immediately be detected and measured with the same ease, speed and precision as changes in declination can be with the ordinary magnetometer.

Several years ago I had already confirmed that this idea worked in practice by preliminary experiments with a device that was admittedly quite crude. I alluded to it in my paper on the Earth’s magnetism and the magnetometer [p. 327 of this volume]. Recently, however, I have had a more refined instrument made and have hung it in the astronomical observatory where the magnetometer with the twenty-five-pound bar used to hang. After what has already been explained, a brief description of the instrument will suffice.

It hangs from two steel wires 17 feet long, or, to be precise, from a single wire whose ends are tied to the instrument below and whose middle runs over two cylinders above. The cylinders keep the two strands a suitable distance apart (about $1\frac{1}{2}$ inches). This arrangement also has the advantage that the two strands naturally have the same tension. The suspension is above the ceiling of the room, and the wires hang freely through a circular opening $3\frac{1}{2}$ inches across in the ceiling. The distance between the wires can be increased or decreased as required at both the top and the bottom. The instrument hanging from the wires itself consists of four main parts. The first, to which the wires are attached, is a horizontal circular disc four inches in diameter, with quarter-degree divisions marked on silver. The second part consists of an alidade that can turn on the graduated rim of the circle, concentric with the circle, and has two verniers that read to individual minutes of arc. A fairly stout rod perpendicular to the plane of the circle is firmly attached to the alidade. Mounted on the rod is a mirror of very exact circular form $1\frac{1}{2}$ inches in diameter. With a telescope 16 feet away, one sees in this mirror the image of part of a horizontal scale marked in individual millimetres and fixed below the telescope. In this way \SourceLeaf{369}every change in the position of the circle can be detected and measured. Small changes are measured directly with the utmost precision by the scale divisions visible in the telescope. Larger ones are measured by combining this with a movement of the alidade and reading the verniers. The third part is the cradle beneath the circle, a double frame through which the fourth component, a strong twenty-five-pound bar magnet, is inserted. This cradle can also turn about the centre of the circle and has two verniers resting on the graduated rim, so that the amount of rotation can be measured to the minute of arc.

If we first set the cradle so that the instrument keeps the same equilibrium position whether the cradle holds the bar magnet or a nonmagnetic body of the same weight, this gives the first or second of the main positions distinguished earlier, according as the bar magnet is in the natural or reversed position. The first offers no particularly important practical application, and the second is useful only on condition that the magnetic directive force is somewhat smaller than the directive force due to the suspension. In the instrument here the ratio of these directive forces is now nearly 10 to 11. The resulting directive force is thus only one tenth of the magnetic directive force. We have here, then, the analogue of an astatic magnetic needle. Any external force that disturbs the direction of a simple needle has ten times as great an effect here as it would with suspension from a single thread, and, as is easy to see, in the opposite direction. This provides, among other things, a solution to a problem that people had repeatedly tackled before without success, namely how to show the daily and hourly variations in magnetic declination on an enlarged scale. Repeated simultaneous observations of this kind with this instrument and with the magnetometer of the magnetic observatory have always given the most satisfactory results. Meanwhile this application is now losing some of its importance, since ordinary magnetometers already show the smallest changes with all the precision one could wish for, so there is no longer a need for magnification.

This and other applications with the bar mounted in the reversed position, to \SourceLeaf{370}which I shall return later, are of secondary importance. Far more important is the use of the instrument in the third or \textit{transverse} position for measuring intensity. If, starting from the natural position, one turns the cradle to move the bar magnet out of the magnetic meridian, the entire instrument must turn back through a certain angle corresponding to the ratio of the two directive forces to reach equilibrium. The difference between the two angles will be the deviation of the bar magnet from the magnetic meridian in the equilibrium position. One can readily arrange for this deviation to be nearly 90 degrees, thereby gaining the advantages mentioned earlier. The instrument is then particularly well suited to observing the \textit{changes} in intensity, which show themselves directly in its change of position. It is in the nature of the matter that several circumstances must be taken into account for changes that occur only after a longer time. In particular, they require us to use (known) suitable means from time to time to examine whether and by how much the strength of the magnetism in the bar has changed. Temperature changes also have to be considered, both insofar as they affect this strength and insofar as they affect the distance between the suspension wires and their length, and hence the directive force due to the suspension. For irregular changes in intensity that vary over short intervals, however, the instrument does exactly what the magnetometer does for similar changes in declination. The method of observation is also exactly the same with both instruments. The changes in intensity are obtained initially in scale divisions, which can easily be converted into fractions of the intensity itself. With the instrument as it is now, one scale division corresponds to the 22000\textsuperscript{th} part of the total intensity.

The experience gained with the instrument so far, though it covers only a short time and is based on relatively few observations, already points to some results of no small importance.

First, the observations so far suggest regular changes depending on the time of day. These may admittedly be mingled with irregular ones just as often as in the case of declination, and separating them with certainty \SourceLeaf{371}will require observations continued for years. If, on the strength of the little experience available so far, I may venture to offer a conjecture rather than a result, the regular course seems to be that the intensity decreases in the morning, already reaches its minimum one or two hours before noon, and increases again from then on. To get at least a preliminary indication of the quantitative relationship, in August 1837 I recorded the position on 30 days, at 10 in the morning and 3 in the afternoon. On 26 days the intensity was greater in the afternoon than in the morning, and on only 4 days was it smaller. The mean difference was 39 scale divisions, or a little more than the 600\textsuperscript{th} part of the total intensity. On most of those days the reading was also recorded at 9 in the morning. Of 28 days there were 23 on which the intensity at that hour was still greater than it was an hour later, and on only 5 days was the reverse the case: the mean difference here, however, was only $11\frac{1}{2}$ scale divisions, or a little more than the 2000\textsuperscript{th} part of the total intensity.

\expandafter\def\csname HistoricalNoteText371\endcsname{Similar observations were later made with both instruments during the scheduled observation session of 13---14 November, with the same result.}
Second, several particularly extensive series of observations confirm that irregular disturbances, sometimes very considerable and varying over short intervals, occur no less often in intensity than in declination, something that in any case could not be doubted on the grounds of analogy alone. On three occasions readings were taken without interruption for a considerable time with this intensity instrument and simultaneously with the magnetometer of the magnetic observatory. The first was on 15 July, from 6 in the morning to 6 in the evening. Then came the regular magnetic observation session of 29---30 July, and finally the special session from 31 August to 1 September, 24 hours on each of these last two occasions. Readings were always taken every 5 minutes. Graphical representations of the two sessions, with curves drawn for changes in both intensity and declination, bring this out clearly. The two motions, as goes without saying, bear not the slightest resemblance to each other. Yet it is very noticeable that where declination was strongly disturbed, strong disturbances usually occurred in the intensity as well\HistoricalNote{371}.
\expandafter\def\csname HistoricalNoteText376\endcsname{The first public mention of these experiments is in the Gött. gel. Anz. 1834, Aug. 9. See \textsc{Schumacher}’s Jahrbuch for 1836, [p. 339 of this volume].}
\expandafter\def\csname HistoricalNoteText376b\endcsname{Gött. gel. Anz. 1835, March 7. \textsc{Schumacher}’s Jahrbuch for 1836, [p. 341 of this volume.]}
\par
\SourceLeaf{372}Plotting the changes in declination and intensity as two separate curves, incidentally, gives a far less vivid picture of how the disturbances unfold than combining them in a single curve. The clearest way to see what matters here is as follows. A single straight line can give a complete picture of the Earth’s magnetic force (its horizontal component, as is always tacitly understood) at any instant. Its length is proportional to the intensity, and it makes an angle equal to the declination with a fixed straight line. To represent the force at several successive instants, we keep the starting point of the different straight lines fixed. Only their endpoints then need to be considered. These can be labelled with the corresponding times and joined by a line. The straight radii themselves are omitted entirely, and even the common starting point will always lie far outside the drawing if the representation is on any reasonably suitable scale. This way of proceeding also gives us a new way of looking at changes of this kind in the two magnetic elements. They are in fact simply the two horizontal components of the always comparatively very small disturbing force to which the mean magnetic force of the Earth is subject at any instant. The disturbing force is resolved in two directions, one \textit{in} the magnetic meridian and the other perpendicular to it. The magnetometer gives the second component directly, and the new instrument gives the first. The two need only be brought to a common scale before plotting.

One further point about using this otherwise very vivid way of representing the changes needs mention here. If the disturbances are frequent and change rapidly, it is difficult to plot the course of an entire day in one drawing without confusion, as the curve develops a great many entanglements. In that case, shorter intervals have to be plotted separately.

If we compare what the new instrument and the magnetometer can do, we find that for \textit{some} purposes they must complement one another, while for others they are equally applicable. Only the \SourceLeaf{373}magnetometer can determine absolute declination. Changes in declination, especially those that vary rapidly, can be tracked with both instruments. Both can determine absolute intensity, although using the magnetometer is somewhat less complicated than using the new instrument alone would be. The magnetometer by itself, however, can give only the mean intensity over a certain interval. It entirely misses rapid variations in intensity, which the new instrument shows most satisfactorily. For all other uses, such as comparing bar magnets with one another in magnetic strength, or for galvanometry and telegraphy in conjunction with a multiplier, both instruments are equally useful. For the last two purposes the new instrument even has a considerable advantage. As already noted, we can make it as nearly astatic as we like.

A few examples may be given here of the instrument’s sensitivity as a galvanometer. The multiplier surrounding the bar magnet contains 610 turns of silk-covered copper wire. In this coil alone a galvanic current has to pass through more than 6000 feet of wire. The length increases to 13000 feet when the current also goes to the physics cabinet. Other instruments are usually added to the circuit, however, so that in some experiments the total length of wire is 40000 feet, or almost two Meilen. It should also be noted that by far the greater part of this wire is much thinner. As far as its effect on the strength of the current is concerned, this length is equivalent to about 8 Meilen of wire of the same thickness as the wire connecting the astronomical observatory to the physics cabinet. Even with such a long circuit, the weakest galvanic forces give the heavy bar magnet a deflection that can be detected and is large enough for precise measurements. This holds, for example, for thermogalvanism, which some physicists mistakenly believe cannot pass through a very long circuit. With the arrangements here, and using a thermogalvanic device of a particular design, touching the junction with a finger is enough to produce that effect.

\SourceLeaf{374}Applying the instrument to ordinary frictional electricity suggests another interesting observation. As is well known, \textsc{Colladon} discovered that when passed through a multiplier it deflects a magnetic needle in much the same way as a hydrogalvanically excited current. His experiments were at first doubted and later confirmed by \textsc{Faraday}. The latter was the first to make clear that even in a very powerful electrical battery no more electricity is developed than very weak means of hydrogalvanic excitation drive through a wire of moderate length in a few seconds. Experiments with the instruments here had also confirmed several years earlier both the reality and the small magnitude of the electromagnetic effect of electricity generated by a machine. It seemed worth repeating these experiments with the much more sensitive new instrument. Here the procedure was to connect only the conductor and the rubbing pad of an electrical machine in the physics cabinet to the ends of the wire circuit leading to the astronomical observatory, instead of discharging a Leyden jar or a battery of jars through the wire circuit (as \textsc{Colladon} and \textsc{Faraday} had done). The circuit measured 13000 feet including the multiplier. The machine was turned continuously at a uniform speed. At one revolution per second this kept the twenty-five-pound bar magnet in the new instrument at the astronomical observatory deflected by 144 scale divisions (somewhat over 50 minutes of arc), positively or negatively according to the direction in which the electricity passed through the multiplier. The experiments were as regular as one could wish. What seems especially remarkable is that the electromagnetic effect stayed the same even when the circuit was made a whole Meile long by adding other instruments. This might appear to be a substantial difference from currents produced hydrogalvanically, thermogalvanically or by induction. The intensity of those currents, as shown by the magnitude of their electromagnetic effects, always decreases as the conducting path is lengthened. I find in this simply a striking confirmation of the theory according to which the unequal intensities of two galvanic currents, as revealed by their unequal electromagnetic effects, amount only to unequal quantities of electricity passing through each cross section of the conducting path in a given time.
\SourceLeaf{375}With the other methods of producing a current, the greater the resistance the longer circuit presents to the current, the less electricity a given electromotive force develops in a given time. In our experiment, by contrast, the amount of electricity set in motion depends only on the operation of the machine, and \textit{all} the electricity that jumps to the conductor in the form of sparks must traverse the entire circuit, short or long, to balance the opposite electricity from the rubbing pad.

To show the advantage of the new instrument over the magnetometer in electromagnetic telegraphy as well, it will first be necessary to look a little more closely at \textit{how} galvanic currents produce telegraphic signals.

\expandafter\def\csname HistoricalNoteText375\endcsname{According to a note communicated to me by Mr. von \textsc{Humboldt}, ten years earlier \textsc{Bétancourt} had already laid a wire circuit from Aranjuez to Madrid, through which the discharge of a Leyden jar was to serve as a telegraphic signal. Although the details of its outcome do not seem to be known, there can be no doubt that an experiment of this kind would succeed if carried out appropriately. Such a method would still presumably be limited to signalling yes or no in reply to one or a couple of questions agreed on in advance.}
As soon as it was known that the effects of a \textsc{Volta}ic pile could propagate through a very long circuit, it was natural to think of using these forces of nature for telegraphy. Nearly 30 years ago\HistoricalNote{375}, when only a small part of the effects of galvanism was known, \textsc{Sömmerring} already proposed using the release of gas for that purpose. The magnetic effects of galvanic currents, which became known only later, are far better suited to composite signals. Yet it is striking that quite a number of years passed after \textsc{Oersted}’s discovery before anyone seems to have thought of using them in this way. Of course, no well-founded judgment about their large-scale application is possible without precise quantitative knowledge of how galvanic currents are weakened by the length and properties of the conducting wires. Before \textsc{Ohm} and \textsc{Fechner}, ideas on this subject were very incomplete and mistaken. In 1833 a wire connection was made between our astronomical observatory and the physics cabinet, chiefly so that similar investigations into the law governing the strength of galvanic currents under different conditions could be conducted on a large scale. The credit \SourceLeaf{376}for the very difficult execution of this ambitious installation belongs to Professor \textsc{Weber} alone. From the very beginning the circuit was often used for simple telegraphic signals to compare the clocks daily and, experimentally, for composite signals as well. And the possibility of signalling letters, words and whole phrases. was thus already an evident fact\HistoricalNote{376}. These experiments used a hydrogalvanic current excited by weak means alone, namely a single or a double pair of plates and water without added acid. I will pass over a detailed description of the method used at the time, since I later replaced it with a completely different one. That method still had a drawback. With our single circuit and with instruments for which such experiments were only a secondary purpose, no more than two letters could be signalled in one minute. Even with modified apparatus designed expressly for telegraphy, this speed could hardly have been increased much as long as only a single circuit was used. The length of the circuit and the distance between its endpoints clearly have nothing to do with this speed. With a multiple circuit the speed could have been increased greatly, but there was insufficient reason to install one here. The outcome itself could not be in any doubt, and the actual scientific benefit of such a multiple circuit would have been out of proportion to its considerable cost.

The theory of the laws of induction, on the other hand, led me to a completely different method. With this method a single circuit has already served for much faster telegraphy for more than two years, with complete success. I may be all the more justified in dwelling on it a little longer, since I have yet to publish any details of it.

Several years ago I described elsewhere the device I call an inductor{\renewcommand{\thefootnote}{**)}\HistoricalNote{376b}}. I should point out, however, that in place of the inductor with 1050 \SourceLeaf{377}turns described in the first notice, and the later one increased to 3537 turns, one with 7000 turns is now in use, in which the wire alone is more than 7000 feet long. A very simple operation with this inductor sends two strong galvanic currents in opposite directions through the conducting wire in quick succession, each lasting only an extremely short time. One slips the inductor quickly off the double bar magnet over which it was initially placed and immediately returns it to its former position without turning it around. Wherever a magnetic needle surrounded by a multiplier is placed in the circuit, the effect of the two currents is to impart a very high velocity to it for an instant and then immediately cancel that velocity completely. Thus the needle makes a very quick but small movement, to the right or the left as desired, and at once comes to a complete standstill again.

It is obvious that these jerks can be varied and combined in many ways and used to signal letters. Some practice will be needed both to give the signals as quickly and precisely as possible and to read them easily and quickly. Even without acquiring any special practice, however, one can readily signal about seven letters in one minute, as repeated experience has shown. If special \textit{mechanical} devices were provided for the operation, the speed and precision could undoubtedly be increased considerably.

The new instrument has a considerable advantage over the magnetometer in this very kind of telegraphy, for the following reasons. The two opposite impulses that make up one simple signal are, in \textit{strength}, \textit{exactly equal}, so the second cancels exactly as much velocity as the first has produced. Even so, the needle cannot be absolutely at rest between signals, since it can be at rest only in its natural equilibrium position. Even if it is in this position \textit{before} a signal, it is displaced somewhat, however slightly, \textit{by} the signal itself, and the directive force acting on the needle then tends to bring it back. Even though one signal can thus give rise only to an extremely \SourceLeaf{378}slight movement, the accumulation of a large number of signals can result in a substantial displacement from the natural equilibrium position. There may consequently be enough motion between signals to make them somewhat less distinct. As a little thought shows, other things being equal this disturbance is more detrimental when the needle whose jerks carry the signals has a short swing duration than when it has a long one. The disturbance is therefore more marked in the magnetometer of the magnetic observatory than in the one with a twenty-five-pound needle that used to hang in the astronomical observatory. It is less marked still in the new instrument that has now taken its place, if the bar magnet in its reversed position is set up as a nearly astatic needle. Indeed, even after it has moved substantially away from its equilibrium position, the comparatively weak directive force drawing it back does not set it in motion in a way that seriously interferes with the signals. At the same time the current in the multiplier acts on it just as strongly, so that the jerks are just as large as if the needle belonged to an ordinary magnetometer.

A special device I recently had made is exceptionally useful in dealing with the drawbacks and inconveniences of untimely oscillations, both in this kind of telegraphy and in many other uses of magnetic instruments. I call this device a \textit{damper}, because it completely stops oscillations in a very short time that would otherwise continue for many hours, diminishing only very slowly. The damper, made for the time being only for the magnetometer of the magnetic observatory, does this exceptionally well. Even the largest oscillations die out completely in a few minutes. A similar device can be fitted to any oscillating needle, in a magnetometer or in the new instrument at issue here. It will have to form an essential part of every instrument intended to be used seriously for telegraphy by the method described here. A fuller explanation of this device would take us too far from the subject at hand.

\SourceLeaf{379}The new instrument has been given no particular name so far. From its most important use, one might call it an intensity meter. Since it serves for just as wide a range of precise magnetic measurements as the magnetometer, it might equally well lay claim to \textit{the same} name. The chief difference is that the new instrument is suspended from \textit{two} threads, thereby obtaining a \textit{new} directive force with which the magnetic directive force is commensurable. The other differences, in particular how the mirror is mounted and how the rotation of the individual parts relative to one another is measured, are necessary conditions that follow from the aims themselves. One might therefore call the new instrument a \textit{bifilar} or \textit{bipensil magnetometer} to distinguish it from the older instrument, the simple or unifilar magnetometer. I may venture to express my conviction that its wider adoption, especially its use in scheduled observations alongside the simple magnetometer at several places far apart, will soon be followed by substantial advances in our knowledge of the astonishing disturbances of the Earth’s magnetism.

